\appendices
\section{系统环境配置与部署说明}

\subsection{后端环境搭建}

本系统后端基于 Node.js 20 LTS + Express 4 开发，以下为本地环境的完整搭建步骤。

\textbf{（1）安装 Node.js}

从 \texttt{https://nodejs.org/} 下载并安装 Node.js 20 LTS 版本，安装完成后验证：

\begin{lstlisting}[language=bash, caption={验证 Node.js 安装}]
node -v   # 应输出 v20.x.x
npm  -v   # 应输出 10.x.x
\end{lstlisting}

\textbf{（2）安装后端依赖}

\begin{lstlisting}[language=bash, caption={安装后端依赖包}]
cd backend
npm install
\end{lstlisting}

\texttt{package.json} 核心依赖清单如表~\ref{tab:A-1}~所示。

\begin{table}[H]
    \centering
    \caption{后端核心依赖包清单}
    \label{tab:A-1}
    \begin{tabular}{lll}
        \hline
        包名 & 版本 & 说明 \\
        \hline
        express & \^{}4.18.2 & Web 框架 \\
        sequelize & \^{}6.35.2 & ORM 框架 \\
        mysql2 & \^{}3.6.5 & MySQL 驱动 \\
        jsonwebtoken & \^{}9.0.2 & JWT 签发与验证 \\
        bcryptjs & \^{}2.4.3 & 密码哈希 \\
        cors & \^{}2.8.5 & 跨域请求处理 \\
        dotenv & \^{}16.4.1 & 环境变量加载 \\
        \hline
    \end{tabular}
\end{table}

\textbf{（3）配置环境变量}

在 \texttt{backend/} 目录下创建 \texttt{.env} 文件（参照 \texttt{.env.example}）：

\begin{lstlisting}[language=bash, caption={后端 .env 配置示例}]
DB_HOST=localhost
DB_PORT=3306
DB_NAME=mental_health
DB_USER=root
DB_PASSWORD=your_db_password
JWT_SECRET=your_jwt_secret_key
JWT_EXPIRES_IN=7d
DEEPSEEK_API_KEY=your_deepseek_api_key
DEEPSEEK_BASE_URL=https://api.deepseek.com
PORT=3000
\end{lstlisting}

\textbf{（4）启动后端服务}

\begin{lstlisting}[language=bash, caption={启动后端开发服务器}]
npm run dev
\end{lstlisting}

后端服务默认运行于 \texttt{http://localhost:3000}。


\subsection{前端环境搭建}

本系统前端基于 Vue~3 + Vite + Pinia 开发。

\textbf{（1）安装前端依赖}

\begin{lstlisting}[language=bash, caption={安装前端依赖包}]
cd frontend
npm install
\end{lstlisting}

\texttt{package.json} 核心依赖清单如表~\ref{tab:A-2}~所示。

\begin{table}[H]
    \centering
    \caption{前端核心依赖包清单}
    \label{tab:A-2}
    \begin{tabular}{lll}
        \hline
        包名 & 版本 & 说明 \\
        \hline
        vue & \^{}3.4.x & 前端核心框架 \\
        vue-router & \^{}4.2.x & 单页路由管理 \\
        pinia & \^{}2.1.x & 状态管理 \\
        axios & \^{}1.6.x & HTTP 客户端 \\
        vite & \^{}5.0.x & 构建工具 \\
        \hline
    \end{tabular}
\end{table}

\textbf{（2）配置前端环境变量}

在 \texttt{frontend/} 目录下创建 \texttt{.env} 文件：

\begin{lstlisting}[language=bash, caption={前端 .env 配置}]
VITE_API_BASE_URL=http://localhost:3000
\end{lstlisting}

\textbf{（3）启动前端开发服务器}

\begin{lstlisting}[language=bash, caption={启动前端开发服务器}]
npm run dev
\end{lstlisting}

前端默认运行于 \texttt{http://localhost:5173}。

\textbf{（4）生产环境构建}

\begin{lstlisting}[language=bash, caption={生产构建}]
npm run build
\end{lstlisting}

构建产物输出至 \texttt{dist/} 目录，可由 Nginx 或 Node.js 静态服务托管。


\section{数据库初始化说明}

\subsection{创建数据库}

使用 MySQL 8 客户端执行以下命令创建数据库：

\begin{lstlisting}[language=SQL, caption={创建数据库}]
CREATE DATABASE mental_health
  CHARACTER SET utf8mb4
  COLLATE utf8mb4_unicode_ci;
\end{lstlisting}

\subsection{执行 Migration 脚本}

数据库表结构通过 \texttt{database/migrations/} 目录下按序号排列的 SQL 脚本管理。按编号顺序在 \texttt{mental\_health} 数据库中依次执行各脚本：

\begin{lstlisting}[language=bash, caption={执行 Migration 脚本}]
mysql -u root -p mental_health < database/migrations/001_init.sql
\end{lstlisting}

所有 Migration 脚本均采用幂等写法（\texttt{CREATE TABLE IF NOT EXISTS} 等），可重复执行而不产生副作用。执行完成后，启动后端服务，Sequelize 将在首次运行时通过模型同步完成剩余的表结构校验。

